\documentclass[11pt,a4paper]{article}
\usepackage[utf8]{inputenc}
\usepackage[T1]{fontenc}
\usepackage[margin=2.2cm]{geometry}
\usepackage{amsmath,amssymb}
\usepackage{enumitem}
\usepackage{hyperref}

\setlength{\parindent}{0pt}
\setlength{\parskip}{0.75em}

\title{MA087G Mathematical Statistics\\Beginner-Friendly Full Solution Guide\\2025-06-07}
\author{}
\date{}

\begin{document}
\maketitle

\section*{Source notes before solving}
The official exam PDF was image-only. Normal text extraction found no text, but the rendered page images were readable.

Per request, this folder does not include an extraction log or mistake bank. The guide below uses the updated beginner-friendly style.

\section*{Uppgift 1}

\subsection*{What this question is really asking}
There is a jar with:
\[
10 \text{ red balls} \quad \text{and} \quad 10 \text{ white balls}.
\]
We choose 5 balls without replacement. ``Without replacement'' means we do not put a ball back after drawing it.

Event \(A\): all 5 balls are red.

Event \(B\): we get at least one ball of each color. So \(B\) means at least one red and at least one white.

We must calculate:
\[
P(A),\quad P(B),\quad P(A\cup B),\quad P(A\mid B^c).
\]

\subsection*{How to recognize this type of question}
The clues are:
\begin{itemize}[leftmargin=*]
    \item We draw several balls from a finite group.
    \item It says ``utan återläggning'', meaning without replacement.
    \item The order of the balls does not matter.
\end{itemize}

This is a combinations/hypergeometric problem. We count possible groups of 5 balls.

\subsection*{Tiny theory needed}
When order does not matter, use combinations:
\[
\binom{n}{k}.
\]
This means ``choose \(k\) objects from \(n\) objects''.

The total number of possible draws is:
\[
\binom{20}{5},
\]
because there are 20 balls total and we choose 5.

The probability is:
\[
\frac{\text{number of favorable groups}}{\text{number of all possible groups}}.
\]

\subsection*{Formula and what everything means}
Combination formula:
\[
\binom{n}{k}=\frac{n!}{k!(n-k)!}.
\]
Here:
\begin{itemize}[leftmargin=*]
    \item \(n\) is how many objects we choose from.
    \item \(k\) is how many objects we choose.
    \item The exclamation mark means factorial, for example \(5!=5\cdot4\cdot3\cdot2\cdot1\).
\end{itemize}

Conditional probability:
\[
P(A\mid B^c)=\frac{P(A\cap B^c)}{P(B^c)}.
\]
Here \(B^c\) means ``not \(B\)''.

\subsection*{Step-by-step solution}
The total number of ways to choose 5 balls from 20 is:
\[
\binom{20}{5}=15504.
\]

\textbf{Part a: Calculate \(P(A)\).}

Event \(A\) means all 5 balls are red. There are 10 red balls, so the number of favorable choices is:
\[
\binom{10}{5}=252.
\]
Therefore:
\begin{align*}
P(A)
&=\frac{\binom{10}{5}}{\binom{20}{5}}\\
&=\frac{252}{15504}\\
&=\frac{21}{1292}\\
&\approx 0.0163.
\end{align*}
So all red is quite unlikely.

\textbf{Part b: Calculate \(P(B)\).}

Event \(B\) means at least one of each color.

Instead of counting all mixed cases directly, use the complement. Not \(B\), written \(B^c\), means we do not have both colors. That can only happen in two ways:
\[
\text{all red} \quad \text{or} \quad \text{all white}.
\]
The probability of all red is the same as the probability of all white:
\[
\frac{\binom{10}{5}}{\binom{20}{5}}.
\]
So:
\begin{align*}
P(B^c)
&=P(\text{all red})+P(\text{all white})\\
&=2\cdot \frac{\binom{10}{5}}{\binom{20}{5}}\\
&=2\cdot \frac{252}{15504}\\
&=\frac{504}{15504}.
\end{align*}
Then:
\begin{align*}
P(B)
&=1-P(B^c)\\
&=1-\frac{504}{15504}\\
&=\frac{15000}{15504}\\
&=\frac{625}{646}\\
&\approx 0.9675.
\end{align*}

\textbf{Part c: Calculate \(P(A\cup B)\).}

The symbol \(\cup\) means ``or''. So \(A\cup B\) means:
\[
A \text{ happens or } B \text{ happens}.
\]

Event \(A\) is all red. Event \(B\) is at least one red and at least one white.

These cannot happen at the same time. If all balls are red, then we do not have a white ball. So:
\[
A\cap B=\varnothing.
\]
That means:
\[
P(A\cap B)=0.
\]
Use:
\begin{align*}
P(A\cup B)
&=P(A)+P(B)-P(A\cap B)\\
&=P(A)+P(B)-0.
\end{align*}
Insert the values:
\begin{align*}
P(A\cup B)
&=\frac{21}{1292}+\frac{625}{646}\\
&=\frac{21}{1292}+\frac{1250}{1292}\\
&=\frac{1271}{1292}\\
&\approx 0.9837.
\end{align*}

Another way to see it: \(A\cup B\) is everything except all white.

\textbf{Part d: Calculate \(P(A\mid B^c)\).}

We want the probability of \(A\), given that \(B\) did not happen.

Remember:
\[
B^c=\text{all red or all white}.
\]
Event \(A\) is all red. If \(A\) happens, then \(B^c\) definitely happens. So:
\[
A\cap B^c=A.
\]
Use the conditional probability formula:
\begin{align*}
P(A\mid B^c)
&=\frac{P(A\cap B^c)}{P(B^c)}\\
&=\frac{P(A)}{P(B^c)}.
\end{align*}
Now insert values:
\begin{align*}
P(A\mid B^c)
&=\frac{\frac{252}{15504}}{\frac{504}{15504}}\\
&=\frac{252}{504}\\
&=\frac{1}{2}.
\end{align*}

\subsection*{Answer}
\[
\boxed{P(A)=\frac{21}{1292}\approx 0.0163}
\]
\[
\boxed{P(B)=\frac{625}{646}\approx 0.9675}
\]
\[
\boxed{P(A\cup B)=\frac{1271}{1292}\approx 0.9837}
\]
\[
\boxed{P(A\mid B^c)=\frac{1}{2}}
\]

\subsection*{What to memorize for the exam}
If the problem says ``without replacement'' and order does not matter, think combinations:
\[
\frac{\text{good combinations}}{\text{all combinations}}.
\]

\subsection*{Common trap}
Do not use a binomial formula with \(p=1/2\) here. The draws are without replacement, so the probability changes after each draw. Combinations handle this cleanly.

\subsection*{Repeated exam pattern}
This is a new pattern in the processed folders so far: hypergeometric/combinations without replacement. It still connects to the probability-event thinking from 2026-01 Uppgift 1 and 2025-08 Uppgift 1.

\section*{Uppgift 2}

\subsection*{What this question is really asking}
We are given a function:
\[
f(x)=
\begin{cases}
\dfrac{a}{\sqrt{x}}, & 0<x\le 4,\\
0, & \text{otherwise}.
\end{cases}
\]
We need to make it into a density function, then calculate the CDF, a probability, expected value, and standard deviation.

\subsection*{How to recognize this type of question}
The clues are:
\begin{itemize}[leftmargin=*]
    \item It says ``täthetsfunktion'', meaning density function.
    \item It asks for \(F_X\), the distribution function.
    \item It asks for \(P(X\le 1)\), expected value, and standard deviation.
\end{itemize}

So we use integrals.

\subsection*{Tiny theory needed}
A density must have total area \(1\):
\[
\int f(x)\,dx=1.
\]

The CDF is:
\[
F_X(x)=P(X\le x).
\]
For a continuous variable, we find it by integrating the density up to \(x\).

\subsection*{Formula and what everything means}
Density condition:
\[
\int_0^4 \frac{a}{\sqrt{x}}\,dx=1.
\]
Here:
\begin{itemize}[leftmargin=*]
    \item \(a\) is the unknown constant.
    \item \(\sqrt{x}=x^{1/2}\).
    \item \(\frac{1}{\sqrt{x}}=x^{-1/2}\).
\end{itemize}

Expected value:
\[
E(X)=\int x f(x)\,dx.
\]
Variance:
\[
V(X)=E(X^2)-[E(X)]^2.
\]
Standard deviation:
\[
D(X)=\sqrt{V(X)}.
\]

\subsection*{Step-by-step solution}
\textbf{Part a: Find \(a\).}

Use the total area condition:
\begin{align*}
\int_0^4 \frac{a}{\sqrt{x}}\,dx&=1.
\end{align*}
Write \(\frac{1}{\sqrt{x}}\) as \(x^{-1/2}\):
\begin{align*}
a\int_0^4 x^{-1/2}\,dx&=1.
\end{align*}
Integrate:
\[
\int x^{-1/2}\,dx=2x^{1/2}=2\sqrt{x}.
\]
So:
\begin{align*}
a[2\sqrt{x}]_0^4&=1\\
a(2\sqrt{4}-2\sqrt{0})&=1\\
a(4-0)&=1\\
4a&=1\\
a&=\frac{1}{4}.
\end{align*}

\textbf{Part b: Find \(F_X(x)\).}

There are three regions:
\begin{itemize}[leftmargin=*]
    \item If \(x\le 0\), then \(F_X(x)=0\).
    \item If \(0<x\le 4\), integrate from \(0\) to \(x\).
    \item If \(x>4\), then \(F_X(x)=1\).
\end{itemize}

For \(0<x\le 4\):
\begin{align*}
F_X(x)
&=\int_0^x \frac{1}{4\sqrt{t}}\,dt\\
&=\frac{1}{4}\int_0^x t^{-1/2}\,dt\\
&=\frac{1}{4}[2\sqrt{t}]_0^x\\
&=\frac{1}{2}\sqrt{x}.
\end{align*}
Therefore:
\[
F_X(x)=
\begin{cases}
0, & x\le 0,\\[1mm]
\dfrac{\sqrt{x}}{2}, & 0<x\le 4,\\[2mm]
1, & x>4.
\end{cases}
\]

\textbf{Part c: Calculate \(P(X\le 1)\).}

Use the CDF:
\begin{align*}
P(X\le 1)
&=F_X(1)\\
&=\frac{\sqrt{1}}{2}\\
&=\frac{1}{2}.
\end{align*}

\textbf{Part d: Expected value and standard deviation.}

First:
\begin{align*}
E(X)
&=\int_0^4 x\cdot \frac{1}{4\sqrt{x}}\,dx\\
&=\frac{1}{4}\int_0^4 x^{1/2}\,dx\\
&=\frac{1}{4}\left[\frac{2}{3}x^{3/2}\right]_0^4\\
&=\frac{1}{6}\cdot 4^{3/2}\\
&=\frac{1}{6}\cdot 8\\
&=\frac{4}{3}.
\end{align*}

Now calculate \(E(X^2)\):
\begin{align*}
E(X^2)
&=\int_0^4 x^2\cdot \frac{1}{4\sqrt{x}}\,dx\\
&=\frac{1}{4}\int_0^4 x^{3/2}\,dx\\
&=\frac{1}{4}\left[\frac{2}{5}x^{5/2}\right]_0^4\\
&=\frac{1}{10}\cdot 4^{5/2}\\
&=\frac{1}{10}\cdot 32\\
&=\frac{16}{5}.
\end{align*}

Variance:
\begin{align*}
V(X)
&=E(X^2)-[E(X)]^2\\
&=\frac{16}{5}-\left(\frac{4}{3}\right)^2\\
&=\frac{16}{5}-\frac{16}{9}\\
&=\frac{144}{45}-\frac{80}{45}\\
&=\frac{64}{45}.
\end{align*}

Standard deviation:
\begin{align*}
D(X)
&=\sqrt{\frac{64}{45}}\\
&=\frac{8}{3\sqrt{5}}\\
&\approx 1.19.
\end{align*}

\subsection*{Answer}
\[
\boxed{a=\frac{1}{4}}
\]
\[
\boxed{
F_X(x)=
\begin{cases}
0, & x\le 0,\\
\frac{\sqrt{x}}{2}, & 0<x\le 4,\\
1, & x>4
\end{cases}}
\]
\[
\boxed{P(X\le 1)=\frac{1}{2}}
\]
\[
\boxed{E(X)=\frac{4}{3}},\qquad
\boxed{D(X)=\sqrt{\frac{64}{45}}\approx 1.19}
\]

\subsection*{What to memorize for the exam}
For density questions, first normalize, then build the CDF, then calculate moments:
\[
\int f(x)\,dx=1,\qquad F_X(x)=\int f(t)\,dt.
\]

\subsection*{Common trap}
Do not say \(E(X)=1\) just because \(P(X\le 1)=1/2\). Probability and expected value are different things.

\subsection*{Repeated exam pattern}
This is the same pattern as 2026-01 Uppgift 2 and 2025-08 Uppgift 2: continuous density normalization, CDF, probability, expected value, and standard deviation.

\section*{Uppgift 3}

\subsection*{What this question is really asking}
We have:
\[
X\sim \operatorname{Exp}(2),\qquad Y\sim \operatorname{Exp}(3),
\]
and \(X,Y\) are independent.

We must calculate one exponential probability, mean and variance of \(3X+Y\), the CDF of \(4X+1\), and a probability involving \(\max(X,Y)\).

\subsection*{How to recognize this type of question}
The clues are:
\begin{itemize}[leftmargin=*]
    \item It says \(\operatorname{Exp}(2)\) and \(\operatorname{Exp}(3)\).
    \item It says independent.
    \item It asks about linear combinations and \(\max\).
\end{itemize}

\subsection*{Tiny theory needed}
If \(X\sim \operatorname{Exp}(\lambda)\), then:
\[
P(X>t)=e^{-\lambda t},
\]
\[
F_X(t)=P(X\le t)=1-e^{-\lambda t},\quad t\ge 0,
\]
\[
E(X)=\frac{1}{\lambda},\qquad V(X)=\frac{1}{\lambda^2}.
\]

\subsection*{Formula and what everything means}
For \(X\sim\operatorname{Exp}(2)\):
\[
P(X>t)=e^{-2t}.
\]
Here \(2\) is the rate.

For a linear combination:
\[
E(3X+Y)=3E(X)+E(Y).
\]
If independent:
\[
V(3X+Y)=3^2V(X)+V(Y).
\]

\subsection*{Step-by-step solution}
\textbf{Part a: Calculate \(P(1\le X<3)\).}

For exponential variables, interval probability is:
\[
P(1\le X<3)=F_X(3)-F_X(1).
\]
Using survival is even faster:
\[
P(1\le X<3)=P(X\ge 1)-P(X\ge 3).
\]
So:
\begin{align*}
P(1\le X<3)
&=e^{-2\cdot1}-e^{-2\cdot3}\\
&=e^{-2}-e^{-6}\\
&\approx 0.1329.
\end{align*}

\textbf{Part b: Calculate \(E(3X+Y)\) and \(V(3X+Y)\).}

For \(X\sim\operatorname{Exp}(2)\):
\[
E(X)=\frac{1}{2},\qquad V(X)=\frac{1}{2^2}=\frac{1}{4}.
\]
For \(Y\sim\operatorname{Exp}(3)\):
\[
E(Y)=\frac{1}{3},\qquad V(Y)=\frac{1}{3^2}=\frac{1}{9}.
\]

Expected value:
\begin{align*}
E(3X+Y)
&=3E(X)+E(Y)\\
&=3\cdot \frac{1}{2}+\frac{1}{3}\\
&=\frac{3}{2}+\frac{1}{3}\\
&=\frac{9}{6}+\frac{2}{6}\\
&=\frac{11}{6}.
\end{align*}

Variance:
\begin{align*}
V(3X+Y)
&=3^2V(X)+V(Y)\\
&=9\cdot \frac{1}{4}+\frac{1}{9}\\
&=\frac{9}{4}+\frac{1}{9}\\
&=\frac{81}{36}+\frac{4}{36}\\
&=\frac{85}{36}.
\end{align*}

\textbf{Part c: Find the CDF of \(4X+1\).}

Let:
\[
W=4X+1.
\]
We want:
\[
F_W(w)=P(W\le w).
\]
Substitute \(W=4X+1\):
\begin{align*}
F_W(w)
&=P(4X+1\le w).
\end{align*}
Solve the inequality for \(X\):
\begin{align*}
4X+1&\le w\\
4X&\le w-1\\
X&\le \frac{w-1}{4}.
\end{align*}
Since \(X\ge 0\), the smallest possible value of \(W=4X+1\) is \(1\). So if \(w<1\), then:
\[
F_W(w)=0.
\]
If \(w\ge 1\):
\begin{align*}
F_W(w)
&=P\left(X\le \frac{w-1}{4}\right)\\
&=1-e^{-2\cdot (w-1)/4}\\
&=1-e^{-(w-1)/2}.
\end{align*}
So:
\[
F_W(w)=
\begin{cases}
0, & w<1,\\
1-e^{-(w-1)/2}, & w\ge 1.
\end{cases}
\]

\textbf{Part d: Calculate \(P(\max\{X,Y\}\le 1)\).}

The maximum is at most 1 only if both variables are at most 1:
\[
\max\{X,Y\}\le 1
\]
means:
\[
X\le 1 \text{ and } Y\le 1.
\]
Because \(X\) and \(Y\) are independent:
\begin{align*}
P(\max\{X,Y\}\le 1)
&=P(X\le 1)P(Y\le 1).
\end{align*}
Now:
\[
P(X\le 1)=1-e^{-2},
\]
and:
\[
P(Y\le 1)=1-e^{-3}.
\]
Therefore:
\begin{align*}
P(\max\{X,Y\}\le 1)
&=(1-e^{-2})(1-e^{-3})\\
&\approx 0.8216.
\end{align*}

\subsection*{Answer}
\[
\boxed{P(1\le X<3)=e^{-2}-e^{-6}\approx 0.1329}
\]
\[
\boxed{E(3X+Y)=\frac{11}{6}},\qquad
\boxed{V(3X+Y)=\frac{85}{36}}
\]
\[
\boxed{
F_{4X+1}(w)=
\begin{cases}
0, & w<1,\\
1-e^{-(w-1)/2}, & w\ge 1
\end{cases}}
\]
\[
\boxed{P(\max\{X,Y\}\le 1)=(1-e^{-2})(1-e^{-3})\approx 0.8216}
\]

\subsection*{What to memorize for the exam}
For exponential:
\[
F(t)=1-e^{-\lambda t},\qquad P(X>t)=e^{-\lambda t}.
\]
For maximum:
\[
\max(X,Y)\le c
\]
means both \(X\le c\) and \(Y\le c\).

\subsection*{Common trap}
Do not treat \(\max(X,Y)\le 1\) as only one variable being below 1. Both must be below 1.

\subsection*{Repeated exam pattern}
Part b repeats the linear-combination expectation/variance pattern from 2025-08 Uppgift 3c and 2026-01 Uppgift 4b. Part d repeats the maximum pattern from 2025-08 Uppgift 3b.

\section*{Uppgift 4}

\subsection*{What this question is really asking}
This is a maximum likelihood question. We have a density with an unknown parameter \(a\):
\[
f_X(x)=
\begin{cases}
\dfrac{x}{a}e^{-x^2/(2a)}, & x\ge 0,\\
0, & \text{otherwise}.
\end{cases}
\]
We need to derive the ML estimate of \(a\).

\subsection*{How to recognize this type of question}
The clue is:
\[
\text{``Härled ML-skattningen av \(a\)''}.
\]
That means:
\[
\text{write likelihood, take log, differentiate, solve.}
\]

\subsection*{Tiny theory needed}
For independent observations \(x_1,\ldots,x_n\), the likelihood is the product of the densities:
\[
L(a)=\prod_{i=1}^n f_X(x_i).
\]
We take the log because products become sums, which are easier to differentiate.

\subsection*{Formula and what everything means}
Likelihood:
\[
L(a)=\prod_{i=1}^n \frac{x_i}{a}e^{-x_i^2/(2a)}.
\]
Here:
\begin{itemize}[leftmargin=*]
    \item \(x_i\) is one observed value.
    \item \(n\) is the number of observations.
    \item \(a\) is the unknown parameter.
\end{itemize}

\subsection*{Step-by-step solution}
Start with:
\begin{align*}
L(a)
&=\prod_{i=1}^n \frac{x_i}{a}e^{-x_i^2/(2a)}.
\end{align*}
Separate the product:
\begin{align*}
L(a)
&=\left(\prod_{i=1}^n x_i\right)a^{-n}
\exp\left(-\frac{1}{2a}\sum_{i=1}^n x_i^2\right).
\end{align*}

Now take the natural log:
\begin{align*}
\ell(a)
&=\ln L(a)\\
&=\sum_{i=1}^n \ln x_i - n\ln a-\frac{1}{2a}\sum_{i=1}^n x_i^2.
\end{align*}

The first term \(\sum \ln x_i\) does not contain \(a\), so its derivative is 0.

Differentiate:
\begin{align*}
\ell'(a)
&=-\frac{n}{a}+\frac{1}{2a^2}\sum_{i=1}^n x_i^2.
\end{align*}

Set the derivative equal to zero:
\begin{align*}
-\frac{n}{a}+\frac{1}{2a^2}\sum_{i=1}^n x_i^2&=0.
\end{align*}
Move the first term:
\begin{align*}
\frac{1}{2a^2}\sum_{i=1}^n x_i^2&=\frac{n}{a}.
\end{align*}
Multiply both sides by \(2a^2\):
\begin{align*}
\sum_{i=1}^n x_i^2&=2na.
\end{align*}
Solve for \(a\):
\begin{align*}
\hat a_{ML}
&=\frac{1}{2n}\sum_{i=1}^n x_i^2.
\end{align*}

\subsection*{Answer}
\[
\boxed{\hat a_{ML}=\frac{1}{2n}\sum_{i=1}^n x_i^2}
\]

\subsection*{What to memorize for the exam}
For ML:
\[
L \rightarrow \ln L \rightarrow \text{differentiate} \rightarrow \text{set equal to zero} \rightarrow \text{solve}.
\]

\subsection*{Common trap}
Do not forget that \(a\) appears both in \(1/a\) and in the exponential. Both parts affect the derivative.

\subsection*{Repeated exam pattern}
This is the same broad ML pattern as 2026-01 Uppgift 5 and 2025-08 Uppgift 4. The density is new, but the ML workflow is the same.

\section*{Uppgift 5}

\subsection*{What this question is really asking}
We have the same 8 normal observations as in 2025-08:
\[
13.4,\;14.1,\;15.1,\;14.3,\;14.9,\;13.9,\;14.4,\;14.0.
\]
We need two 95\% confidence intervals for the mean \(\mu\):
\begin{itemize}[leftmargin=*]
    \item part a: when \(\sigma=0.5\) is known,
    \item part b: when \(\sigma\) is unknown.
\end{itemize}

\subsection*{How to recognize this type of question}
The clues are:
\begin{itemize}[leftmargin=*]
    \item It says ``konfidensintervall för \(\mu\)'', so we estimate a mean.
    \item Part a gives \(\sigma=0.5\), so use a \(z\)-interval.
    \item Part b says \(\sigma\) is unknown, so use a \(t\)-interval with sample standard deviation \(s\).
\end{itemize}

\subsection*{Tiny theory needed}
Known \(\sigma\):
\[
\bar{x}\pm z_{\alpha/2}\frac{\sigma}{\sqrt{n}}.
\]

Unknown \(\sigma\):
\[
\bar{x}\pm t_{\alpha/2,n-1}\frac{s}{\sqrt{n}}.
\]

\subsection*{Formula and what everything means}
\[
\bar{x}=\frac{\sum x_i}{n}.
\]
\(\bar{x}\) is the sample average.

\[
\frac{\sigma}{\sqrt{n}}
\quad \text{or} \quad
\frac{s}{\sqrt{n}}
\]
is the standard error. It tells us how much the sample average typically varies.

\subsection*{Step-by-step solution}
We are given:
\[
n=8,\qquad \sum x_i=114.1,\qquad \sum x_i^2=1629.45.
\]

First calculate the sample mean:
\begin{align*}
\bar{x}
&=\frac{114.1}{8}\\
&=14.2625.
\end{align*}

\textbf{Part a: 95\% interval when \(\sigma=0.5\) is known.}

For 95\%, use:
\[
z_{0.975}=1.96.
\]
Standard error:
\begin{align*}
\frac{\sigma}{\sqrt{n}}
&=\frac{0.5}{\sqrt{8}}\\
&\approx 0.1768.
\end{align*}
Margin of error:
\begin{align*}
1.96\cdot 0.1768&\approx 0.3465.
\end{align*}
Interval:
\begin{align*}
14.2625\pm 0.3465.
\end{align*}
So:
\[
(13.916,\;14.609).
\]

\textbf{Part b: 95\% interval when \(\sigma\) is unknown.}

Now we need the sample standard deviation \(s\).

First calculate:
\begin{align*}
\sum x_i^2-n\bar{x}^2
&=1629.45-8(14.2625)^2\\
&=1629.45-1627.35125\\
&=2.09875.
\end{align*}
Then:
\begin{align*}
s^2
&=\frac{2.09875}{8-1}\\
&=\frac{2.09875}{7}\\
&\approx 0.2998.
\end{align*}
So:
\[
s=\sqrt{0.2998}\approx 0.5476.
\]

Degrees of freedom:
\[
n-1=7.
\]
For 95\%:
\[
t_{0.975,7}\approx 2.365.
\]
Standard error:
\begin{align*}
\frac{s}{\sqrt{n}}
&=\frac{0.5476}{\sqrt{8}}\\
&\approx 0.1936.
\end{align*}
Margin of error:
\begin{align*}
2.365\cdot 0.1936&\approx 0.458.
\end{align*}
Interval:
\begin{align*}
14.2625\pm 0.458.
\end{align*}
So:
\[
(13.805,\;14.720).
\]

\subsection*{Answer}
\[
\boxed{\text{Known }\sigma=0.5:\ (13.916,\;14.609)}
\]
\[
\boxed{\text{Unknown }\sigma:\ (13.805,\;14.720)}
\]

\subsection*{What to memorize for the exam}
If \(\sigma\) is known, use \(z\). If \(\sigma\) is unknown, calculate \(s\) and use \(t\).

\subsection*{Common trap}
Part a and part b look almost the same, but the critical value changes:
\[
z=1.96 \quad \text{versus} \quad t_{0.975,7}=2.365.
\]

\subsection*{Repeated exam pattern}
Part b is the same confidence-interval-for-mean-with-unknown-\(\sigma\) pattern as 2026-01 Uppgift 6a. The data are also reused from 2025-08 Uppgift 5, but here the interval is for \(\mu\), not \(\sigma\).

\section*{Uppgift 6}

\subsection*{What this question is really asking}
A seller claims the average plant height is:
\[
\mu=70\text{ cm}.
\]
The farmer can reject the delivery if:
\[
\mu<70.
\]
So we test whether the true mean is smaller than 70.

We know:
\[
\sigma=10.
\]
So this is a known-\(\sigma\) \(z\)-test.

\subsection*{How to recognize this type of question}
The clues are:
\begin{itemize}[leftmargin=*]
    \item It asks whether the farmer can reject the delivery.
    \item It says the farmer suspects \(\mu<70\).
    \item It gives \(\sigma=10\), so use \(z\), not \(t\).
    \item The risk is 5\%, so the significance level is \(\alpha=0.05\).
\end{itemize}

\subsection*{Tiny theory needed}
We set up:
\[
H_0:\mu=70
\]
against:
\[
H_1:\mu<70.
\]
This is a left-tailed test because the alternative says ``less than''.

For a left-tailed 5\% test, the critical value is:
\[
-1.645.
\]
We reject \(H_0\) if:
\[
z<-1.645.
\]

\subsection*{Formula and what everything means}
\[
z=\frac{\bar{x}-\mu_0}{\sigma/\sqrt{n}}.
\]
Here:
\begin{itemize}[leftmargin=*]
    \item \(\bar{x}\) is the sample average.
    \item \(\mu_0=70\) is the value in the null hypothesis.
    \item \(\sigma=10\) is the known standard deviation.
    \item \(n\) is sample size.
\end{itemize}

\subsection*{Step-by-step solution}
\textbf{Part a: \(n=100\), \(\bar{x}=68\).}

Hypotheses:
\[
H_0:\mu=70,\qquad H_1:\mu<70.
\]

Calculate the standard error:
\begin{align*}
\frac{\sigma}{\sqrt{n}}
&=\frac{10}{\sqrt{100}}\\
&=\frac{10}{10}\\
&=1.
\end{align*}

Calculate the test statistic:
\begin{align*}
z
&=\frac{68-70}{1}\\
&=-2.
\end{align*}

Compare with the critical value:
\[
-2<-1.645.
\]
So we reject \(H_0\).

In words: the sample mean is low enough that the farmer can reject the delivery at 5\% significance.

\textbf{Part b: \(n=200\), \(\bar{x}=69\).}

Use the same hypotheses:
\[
H_0:\mu=70,\qquad H_1:\mu<70.
\]

Standard error:
\begin{align*}
\frac{10}{\sqrt{200}}
&\approx 0.7071.
\end{align*}

Test statistic:
\begin{align*}
z
&=\frac{69-70}{10/\sqrt{200}}\\
&=\frac{-1}{0.7071}\\
&\approx -1.414.
\end{align*}

Compare with:
\[
-1.645.
\]
Since:
\[
-1.414>-1.645,
\]
we do not reject \(H_0\).

In words: the sample mean is lower than 70, but not low enough to reject at the 5\% level.

\subsection*{Answer}
\[
\boxed{H_0:\mu=70,\qquad H_1:\mu<70}
\]
\[
\boxed{\text{Part a: }z=-2,\text{ reject }H_0.\text{ The farmer can reject the delivery.}}
\]
\[
\boxed{\text{Part b: }z\approx -1.414,\text{ do not reject }H_0.\text{ The farmer cannot reject the delivery.}}
\]

\subsection*{What to memorize for the exam}
Known \(\sigma\) means use:
\[
z=\frac{\bar{x}-\mu_0}{\sigma/\sqrt{n}}.
\]
If the alternative is \(\mu<\mu_0\), it is a left-tailed test.

\subsection*{Common trap}
Do not reject just because \(\bar{x}<70\). It must be far enough below 70 compared with the standard error.

\subsection*{Repeated exam pattern}
This is the same known-\(\sigma\) mean-test pattern as 2025-08 Uppgift 6, but here the test is left-tailed instead of right-tailed.

\section*{Mini-cheatsheet from this exam}

\subsection*{Combinations without replacement}
\[
P=\frac{\text{good combinations}}{\text{all combinations}}.
\]
\[
\binom{n}{k}=\frac{n!}{k!(n-k)!}.
\]

\subsection*{Density and CDF}
\[
\int f(x)\,dx=1.
\]
\[
F_X(x)=P(X\le x)=\int_{-\infty}^{x}f(t)\,dt.
\]
\[
V(X)=E(X^2)-[E(X)]^2.
\]

\subsection*{Exponential distribution}
\[
X\sim\operatorname{Exp}(\lambda):
\quad
F_X(x)=1-e^{-\lambda x},\quad
P(X>x)=e^{-\lambda x}.
\]
\[
E(X)=\frac{1}{\lambda},\qquad V(X)=\frac{1}{\lambda^2}.
\]

\subsection*{Linear combinations}
\[
E(aX+bY)=aE(X)+bE(Y).
\]
If independent:
\[
V(aX+bY)=a^2V(X)+b^2V(Y).
\]

\subsection*{Maximum}
\[
\max(X,Y)\le c
\]
means:
\[
X\le c \text{ and } Y\le c.
\]

\subsection*{Maximum likelihood}
\[
L \rightarrow \ln L \rightarrow \text{differentiate} \rightarrow \text{set equal to zero} \rightarrow \text{solve}.
\]

\subsection*{Confidence interval for mean}
Known \(\sigma\):
\[
\bar{x}\pm z_{\alpha/2}\frac{\sigma}{\sqrt{n}}.
\]
Unknown \(\sigma\):
\[
\bar{x}\pm t_{\alpha/2,n-1}\frac{s}{\sqrt{n}}.
\]

\subsection*{Hypothesis test for mean, known sigma}
\[
z=\frac{\bar{x}-\mu_0}{\sigma/\sqrt{n}}.
\]
Left-tailed 5\% test:
\[
\text{reject if }z<-1.645.
\]

\end{document}
